%======================================================================
% Capítulo 2 — Fundamentação Teórica
% Dissertação MAS-Hunt (ABNT, PT-BR)
% Cita apenas chaves existentes em ../bibliografia.bib e
% ../mashunt/references.bib (declaradas no main.tex via \bibliography).
%======================================================================

\chapter{Fundamentação Teórica}
\label{cap:fundamentacao}

Este capítulo estabelece o arcabouço conceitual sobre o qual o MAS-Hunt (\textit{Multi-Agent System for Threat Hunting}) é construído e avaliado. A exposição parte da prática de caça a ameaças e da detecção comportamental, percorre os fundamentos dos grandes modelos de linguagem e dos agentes de inteligência artificial, examina os sistemas multiagente e — central para a contribuição desta dissertação — a segurança dos próprios agentes, antes de detalhar a infraestrutura de telemetria do ecossistema Windows/Elastic, as técnicas de abuso de binários nativos (\textit{Living-off-the-Land Binaries}, LOLBin) catalogadas pelo MITRE ATT\&CK e, por fim, a injeção adversarial de logs. Sempre que pertinente, os conceitos são amarrados aos cinco mecanismos de governança do MAS-Hunt — M1 (Integridade de Memória), M2 (Validação Cruzada), M3 (Resistência a Injeção), M4 (Monitoramento Comportamental) e M5 (Quarentena) — e às quatro condições experimentais (C1-full, C1-nohardening, C2-naive e C3-Elastic), de modo que a metodologia do Capítulo~\ref{cap:metodologia} e a discussão dos resultados encontrem aqui o seu embasamento.

%======================================================================
\section{Caça a Ameaças e Detecção Comportamental}
\label{sec:cacaameacas}

A caça a ameaças (\textit{threat hunting}) consolidou-se como uma mudança de paradigma na defesa cibernética, deslocando o eixo da postura reativa — que aguarda o disparo de um alerta — para uma busca ativa e iterativa por indícios de comprometimento que tenham escapado dos controles automatizados \cite{academia22}. A premissa operacional é assumir a violação (\textit{assume breach}): parte-se do princípio de que um adversário pode já estar presente no ambiente e, por conseguinte, o objetivo passa a ser reduzir o tempo de permanência do atacante (\textit{dwell time}) e mitigar o dano antes que os objetivos da intrusão sejam alcançados \cite{hillier2022turninghuntedhunterthreat}. Essa urgência é quantificada por relatórios setoriais: o tempo mediano global de permanência do atacante mantém-se na ordem de poucos dias a semanas, intervalo suficiente para movimentação lateral e exfiltração de dados \cite{mandiant2025mtrends}, ao passo que violações envolvendo terceiros respondem por parcela expressiva dos incidentes registrados \cite{verizon2025dbir}.

\subsection{Limitações das defesas convencionais}
\label{subsec:limitacoes}

As arquiteturas tradicionais, ancoradas em detecção por assinatura, são inerentemente reativas: mostram-se eficazes contra ameaças conhecidas, mas incapazes de identificar artefatos novos, polimórficos ou de dia zero, para os quais nenhuma assinatura existe \cite{ucci2019}. A resposta histórica a essa limitação foi a análise comportamental, conduzida por sistemas de detecção de intrusão baseados em host (HIDS) ou em rede (NIDS). Tais sistemas, contudo, operam frequentemente de forma isolada: um HIDS pode observar um processo suspeito e um NIDS pode sinalizar uma conexão de rede anômala, mas, sem mecanismos de correlação, falham em conectar eventos dispersos na narrativa coerente de um ataque coordenado. A caça a ameaças surge justamente para suprir essa lacuna de correlação por meio da investigação dirigida por hipóteses \cite{10713082}.

\subsection{O ciclo da caça dirigida por hipóteses}
\label{subsec:ciclo}

A modalidade mais madura de caça a ameaças é a dirigida por hipóteses, na qual o analista formula uma conjectura sobre comportamento adversário plausível — por exemplo, ``um atacante está utilizando \texttt{certutil.exe} para baixar um estágio secundário'' — e então interroga a telemetria disponível em busca de evidências que a confirmem ou refutem \cite{10713082}. Esse ciclo articula-se em etapas de observação, formulação de hipótese, coleta de evidências e classificação. A automação dessas etapas por meio de aprendizado de máquina e de modelos de linguagem é objeto de pesquisa crescente, em trabalhos que vão da análise de logs \cite{10616474} à detecção de malware em dispositivos de Internet das Coisas mediante redes neurais recorrentes \cite{HADDADPAJOUH201888}, passando por surveys dedicados à aplicação da caça a ameaças contra ransomware \cite{9791234} e por abordagens de aprendizado federado para detecção de ameaças persistentes avançadas (APT) \cite{9931222}. Sistemas como o HuntGPT demonstraram a viabilidade de acoplar detecção de anomalias por aprendizado de máquina a modelos de linguagem dotados de IA explicável, oferecendo ao analista justificativas em linguagem natural \cite{ali2023huntgptintegratingmachinelearningbased}; já propostas como o TIRA exploram a automação do ciclo completo de caça, detecção e perfilamento de atores \cite{10750987}, e estudos sobre a autonomia da defesa cibernética delineiam o horizonte de longo prazo dessa transição \cite{academia21}. A semente conceitual do framework M3 do MAS-Hunt — a cadeia \textsc{Observar} $\rightarrow$ \textsc{Hipotetizar} $\rightarrow$ \textsc{Evidência} $\rightarrow$ \textsc{Validação-Cruzada} $\rightarrow$ \textsc{Classificar} — é uma formalização operacional desse ciclo, projetada de modo a ser, simultaneamente, um protocolo de raciocínio e uma barreira de resistência a manipulação adversarial, conforme detalhado na Seção~\ref{sec:segagentes}.

\subsection{A falácia da taxa-base e seus efeitos sobre o falso positivo}
\label{subsec:taxabase}

Um obstáculo estatístico fundamental atravessa qualquer sistema de detecção que busque eventos raros: a falácia da taxa-base (\textit{base rate fallacy}). Mesmo um classificador de elevada acurácia produz um volume avassalador de falsos positivos quando o fenômeno-alvo é raro na população observada \cite{dolan-gavitt2025}. Formalmente, pelo teorema de Bayes, a probabilidade \textit{a posteriori} de que um alerta corresponda a uma ameaça real é
\begin{equation}
P(A \mid B) = \frac{P(B \mid A)\,P(A)}{P(B)},
\label{eq:bayes}
\end{equation}
em que $P(A)$ é a prevalência da ameaça, $P(B\mid A)$ a sensibilidade do detector e $P(B)$ a probabilidade total de um alerta. Quando $P(A)$ é muito pequeno — o que é típico em telemetria corporativa, na qual a esmagadora maioria dos eventos é benigna — o termo de falsos positivos no denominador domina, e o valor preditivo positivo desaba, ainda que sensibilidade e especificidade sejam altas. Essa realidade tem duas consequências de projeto para o MAS-Hunt. Primeiro, justifica a escolha da métrica F2 (Seção~\ref{subsec:lolbin}) e a comparação explícita da Taxa de Falsos Positivos (FPR) contra as regras padrão do Elastic Security (condição C3-Elastic). Segundo, motiva a validação determinística: em vez de confiar na saída estocástica do modelo, exige-se evidência verificável — por exemplo, o acesso comprovado a um artefato-isca (\textit{canary}) de identificador único e difícil de adivinhar — antes de elevar a confiança em um achado \cite{dolan-gavitt2025}, desacoplando a descoberta flexível, conduzida pela IA, da validação rígida e confiável, conduzida por código não-IA.

%======================================================================
\section{Grandes Modelos de Linguagem e Agentes de IA}
\label{sec:llms}

\subsection{Modelos de linguagem como mecanismo de raciocínio}
\label{subsec:llm}

Os grandes modelos de linguagem (\textit{Large Language Models}, LLMs) são redes neurais de larga escala fundamentadas na arquitetura \textit{Transformer}, cujo mecanismo de autoatenção permite modelar dependências de longo alcance em sequências de texto \cite{vaswani2017attention}. Treinados sobre vastos corpora, exibem capacidades emergentes de compreensão, síntese e raciocínio em linguagem natural que viabilizam, sem treinamento específico de tarefa, a interpretação de telemetria de segurança, a formulação de hipóteses e a geração de regras de detecção. Tais capacidades vêm sendo ampliadas por técnicas de computação em tempo de inferência, como o raciocínio latente recorrente, que aprofundam o processo deliberativo do modelo sem aumentar proporcionalmente o número de parâmetros \cite{arxiv250205171}. Dois atributos dos LLMs são, entretanto, decisivos para o projeto de um sistema de defesa: sua saída é, por construção, estocástica — o que torna indispensável um mecanismo de validação e realimentação por especialistas \cite{academia22} — e seu comportamento é integralmente condicionado pelo contexto textual fornecido, propriedade que abre a superfície de ataque por injeção de instruções discutida na Seção~\ref{sec:segagentes}.

\subsection{De modelos a agentes}
\label{subsec:agentes}

Um agente de inteligência artificial é uma estrutura de software que estende um LLM com a capacidade de perceber um ambiente, raciocinar sobre objetivos e executar ações por meio de ferramentas — consultas a bancos de dados, chamadas a APIs, execução de scripts — em um laço iterativo de percepção, decisão e atuação \cite{academia20, guo2024survey}. É essa autonomia instrumental que distingue um agente de um simples \textit{chatbot}: o agente não apenas responde, mas age sobre o sistema. No MAS-Hunt, os agentes executores interagem diretamente com o Elastic Stack, formulando e refinando hipóteses, consultando o Elasticsearch em busca de atividade anômala, criando regras de detecção pela API do Elastic Security e enriquecendo achados com inteligência externa. A mesma autonomia que confere poder ao agente, contudo, é também a origem de sua fragilidade: a literatura recente demonstra que agentes autônomos podem ser explorados de maneiras que não se aplicam a modelos passivos \cite{academia20, NEURIPS2024_eb113910}, tema central da Seção~\ref{sec:segagentes}. Vale ainda registrar a natureza de uso dual dessas capacidades: os mesmos princípios agênticos que sustentam a defesa habilitam a ofensiva, como evidenciado por sistemas de teste de intrusão autônomos baseados em aprendizado por reforço \cite{moreno2025analysis} e por agentes ofensivos de segurança \cite{academia20}. Qualquer arquitetura defensiva que não modele um adversário igualmente capaz é, por isso, fundamentalmente incompleta.

%======================================================================
\section{Sistemas Multiagente}
\label{sec:mas}

Um sistema multiagente (\textit{Multi-Agent System}, MAS) é composto por múltiplos agentes autônomos que interagem entre si e com o ambiente para alcançar objetivos individuais ou coletivos. Esse paradigma é particularmente adequado à segurança cibernética, domínio em que a defesa requer a coordenação de sensores distribuídos e de componentes de decisão. A coordenação entre agentes é um problema clássico da inteligência artificial distribuída, estudado sob a ótica de organizações, planos e escalonamento \cite{durfee1993organizations}, e revigorado pelos LLMs, que passaram a atuar como mecanismo de raciocínio dos agentes \cite{agashe2023llmcoordination}. A pesquisa contemporânea explora desde a propagação de intenções e raciocínio para coordenação \cite{qiu2024collaborativeintelligencepropagatingintentions}, a cooperação proativa \cite{DBLP:journals/corr/abs-2308-11339} e o aprendizado de coordenação \textit{zero-shot} \cite{pmlr-v202-li23au}, até arquiteturas hierárquicas de agentes de linguagem para coordenação humano-IA em tempo real \cite{liu2023llm}, a redução de computação redundante por execução localmente centralizada \cite{bai2024reducingredundantcomputationmultiagent} e o uso de modelos de fundação multimodais para coordenação distribuída \cite{mahmud2025distributedmultiagentcoordinationusing}.

\subsection{Topologias de orquestração e robustez}
\label{subsec:topologias}

A escolha da topologia de orquestração tem efeito direto sobre a robustez do sistema. Arquiteturas hierárquicas — nas quais uma camada superior coordena camadas subordinadas — exibem a menor degradação de desempenho sob condições adversariais quando comparadas a topologias plenamente distribuídas ou a redes \textit{peer-to-peer} \cite{huang2024resilience}. Esse achado fundamenta diretamente a arquitetura de três camadas do MAS-Hunt — \textbf{Conselho de Governança} (\textit{Board}) $\rightarrow$ \textbf{Gerentes} (\textit{Managers}) $\rightarrow$ \textbf{Trabalhadores/Caçadores} (\textit{Workers/Hunters}). Abordagens correlatas reforçam a importância do projeto estrutural: a geração de robustez por adversários auxiliares evolutivos demonstra como exposições controladas a ataques podem endurecer a coordenação \cite{Yuan_Zhang_Xue_Yin_Chen_Guan_Li_Qian_Yu_2023}, ao passo que protocolos de tolerância bizantina, inspirados em mecanismos de consenso de cadeia de blocos, buscam impedir que um único agente comprometido subverta a decisão coletiva \cite{doe2023emergent}. O Conselho de Governança do MAS-Hunt incorpora esse princípio por meio de um mecanismo de consenso deliberativo entre múltiplos agentes, no qual a aprovação de uma diretriz estratégica exige supermaioria, elevando a barreira para que uma falha ou comprometimento isolado altere o rumo da operação.

%======================================================================
\section{Segurança de Agentes de IA}
\label{sec:segagentes}

A integração de agentes autônomos à segurança cibernética introduz uma classe de riscos qualitativamente distinta: o próprio defensor torna-se alvo. Esta seção sistematiza as três famílias de ataque que o MAS-Hunt foi projetado para resistir e que motivam os mecanismos M1 a M5. Trata-se da contribuição conceitual nuclear desta dissertação, pois desloca a discussão do desempenho de detecção para a resiliência do agente.

\subsection{Envenenamento de memória e de bases de conhecimento}
\label{subsec:envenenamento}

Os componentes cognitivos do agente — sua memória de longo prazo e as bases de conhecimento consultadas por geração aumentada por recuperação (\textit{Retrieval-Augmented Generation}, RAG) — constituem alvo privilegiado. Demonstrou-se que uma base pode ser envenenada com uma taxa de contaminação inferior a $0{,}1\%$ e ainda assim induzir comportamento malicioso direcionado com taxa de sucesso de ataque superior a $80\%$, forçando o agente a recuperar e a agir sobre conteúdo adulterado \cite{NEURIPS2024_eb113910}. A insidiosidade do ataque reside em sua furtividade: a contaminação é estatisticamente diluída e dificilmente detectável por inspeção superficial. O MAS-Hunt responde com o mecanismo \textbf{M1 (Integridade de Memória)}, que submete os resultados de recuperação a validação multifonte e a protocolos de sanitização, e com o \textbf{M5 (Quarentena)}, que isola dados anômalos — por exemplo, registros com carimbos de tempo fabricados — antes que contaminem o raciocínio subsequente.

\subsection{Injeção de prompt direta e indireta}
\label{subsec:injecao}

A injeção de prompt explora o fato de que um LLM não distingue, no nível do texto, entre instruções legítimas do operador e instruções incorporadas aos dados que processa. Na variante indireta, o atacante esconde diretivas maliciosas em conteúdo externo que o agente virá a ler — uma página web, um documento ou, no cenário desta dissertação, um registro de log — de modo que a simples leitura desse conteúdo redirecione o comportamento do agente. Em agentes integrados a ferramentas, esse vetor é especialmente perigoso, pois a instrução injetada pode disparar ações reais sobre o sistema; o benchmark InjecAgent quantifica precisamente essa exposição em agentes que operam com ferramentas \cite{zhan2024injecagent}. O MAS-Hunt confronta esse vetor com o mecanismo \textbf{M3 (Resistência a Injeção)}, operacionalizado pelo framework de análise \textsc{Observar} $\rightarrow$ \textsc{Hipotetizar} $\rightarrow$ \textsc{Evidência} $\rightarrow$ \textsc{Validação-Cruzada} $\rightarrow$ \textsc{Classificar}, cuja disciplina exige que a classificação derive exclusivamente de observáveis verificáveis da telemetria, sinalizando — e não obedecendo — qualquer conteúdo que se apresente como instrução. É justamente a eficácia de M3 que o desenho experimental adversarial — quatro modos de injeção combinados a três níveis de furtividade (\textit{overt}, \textit{encoded}, \textit{steganographic}) — pretende mensurar por meio do \textit{Attack Success Rate} (ASR) e de sua variante não-atribuída (NASR).

\subsection{Amplificação de mau funcionamento e jailbreaks infecciosos}
\label{subsec:amplificacao}

A terceira família ataca não a memória nem a entrada, mas o próprio processo de raciocínio. Na amplificação de mau funcionamento (\textit{malfunction amplification}), um erro menor é deliberadamente induzido e propaga-se em um laço de realimentação de ações inúteis, degradando o agente a ponto de configurar uma negação de serviço; taxas de falha superiores a $80\%$ foram reportadas \cite{zhang2024breaking}. Em sistemas multiagente, o risco escala: os chamados \textit{jailbreaks} infecciosos demonstram que uma única entrada adversarial pode desencadear uma infecção exponencial através de uma população de agentes \cite{gu2024agent}, e a avaliação de deriva de risco de conteúdo (\textit{content risk drift}) mostra como interações adversariais sucessivas podem deslocar gradualmente o comportamento do sistema para fora de suas balizas \cite{liu2024crda}. Embora princípios formais de contenção tenham sido derivados, defesas práticas permanecem escassas — precisamente a lacuna que o MAS-Hunt visa preencher. O mecanismo \textbf{M4 (Monitoramento Comportamental)} observa as sequências de ação e o consumo de recursos dos agentes para detectar e conter comportamento anômalo, enquanto o \textbf{M2 (Validação Cruzada)} encaminha os achados dos caçadores a agentes de validação dedicados, que corroboram a evidência por scripts determinísticos e inteligência externa antes da criação de um caso, impedindo que o erro de um único agente se converta em decisão coletiva. A robustez conferida pela orquestração hierárquica \cite{huang2024resilience} complementa esses mecanismos ao limitar a propagação do comprometimento entre camadas. O conjunto M1–M5 corresponde ao endurecimento que distingue a condição \textbf{C1-full} (MAS-Hunt completo) da \textbf{C1-nohardening} (apenas o framework M3, sem M1/M2/M4/M5) e da \textbf{C2-naive} (LLM de passagem única, sem governança), permitindo isolar, por ablação, a contribuição de cada camada de defesa.

\subsection{Governança e conformidade de sistemas de IA}
\label{subsec:governanca}

A operação de agentes autônomos em ambientes corporativos críticos suscita exigências de governança e conformidade que transcendem o desempenho técnico. Estudos comparativos de \textit{frameworks} de governança, risco e conformidade — COBIT, NIST CSF, ISO~27001 e a norma ISO~42001:2023, específica para sistemas de gestão de IA — evidenciam que os arcabouços tradicionais precisam evoluir para endereçar os riscos próprios dos LLMs \cite{McIntosh2024, mcintosh2024cobit}, e propõem a validação por especialista humano no laço como elemento crítico para a integração segura e conforme. Orientações setoriais sobre o uso do COBIT na governança de sistemas de IA \cite{isaca2025leveraging} convergem nessa direção. No plano nacional, o tratamento de telemetria de \textit{endpoints} — que pode conter dados pessoais — situa-se no marco normativo brasileiro da Lei Geral de Proteção de Dados \cite{LGPD2018} e do Marco Civil da Internet \cite{MarcoCivil2014}, ao passo que a tipificação de delitos informáticos \cite{Dieckmann2012} e o debate sobre o futuro da governança de cibersegurança no país \cite{FGV2023} contextualizam a dimensão regulatória da pesquisa. O Conselho de Governança do MAS-Hunt materializa, no plano arquitetural, essa preocupação: além de orientar a estratégia da caça, ele mantém o especialista humano no laço para a aprovação de ações de contenção, preservando uma postura de \textit{analyst-in-the-loop} apropriada à caça automatizada.

%======================================================================
\section{Elastic Security, Sysmon e Telemetria Windows}
\label{sec:telemetria}

\subsection{O Elastic Stack como terreno de caça}
\label{subsec:elastic}

O MAS-Hunt opera sobre o Elastic Stack, que funciona simultaneamente como ambiente operacional e única fonte de verdade. O \textbf{Elasticsearch} é o repositório central e a base de conhecimento, armazenando e indexando a telemetria; o \textbf{Elastic Security} fornece o motor nativo de detecção e o sistema de gestão de casos; e os \textbf{Elastic Agents}, implantados nos \textit{endpoints} e gerenciados pelo Fleet, coletam dados em tempo real a partir de integrações como Elastic Defend, eventos do Windows, Packetbeat (rede) e osquery (inspeção profunda de sistema). Todas as ações dos agentes do MAS-Hunt são executadas por chamadas autenticadas à API REST do Elastic, o que torna a plataforma plenamente programável. A opção por operar sobre telemetria viva de \textit{endpoints} reais é deliberada: detectores avaliados apenas em traços comportamentais controlados podem não representar a diversidade ruidosa dos hospedeiros reais, em razão de desvio de distribuição, variabilidade ambiental e técnicas de evasão de \textit{sandbox} discutidas na literatura de análise de malware \cite{ucci2019}. Ao fundamentar a análise em dados \textit{in situ}, o MAS-Hunt reduz essa limitação por construção.

\subsection{Sysmon e os identificadores de evento}
\label{subsec:sysmon}

A espinha dorsal da telemetria comportamental no Windows é o System Monitor (Sysmon), serviço que registra, no log de eventos do sistema, atividades de granularidade fina identificadas por códigos de evento (\textit{Event IDs}, EID). Para a caça a abuso de binários nativos, o evento mais informativo é o \textbf{EID~1 (criação de processo)}, que captura a imagem executada, a linha de comando completa, o processo pai e o usuário responsável — registrando, portanto, a cadeia de filiação de processos que revela o encadeamento de um ataque. Outros eventos relevantes incluem o EID~3 (conexão de rede), o EID~10 (acesso a processo, útil para detectar despejo de credenciais a partir do \texttt{lsass.exe}) e os EID~12, 13 e 14 (criação e modificação de registro), frequentemente associados a mecanismos de persistência. Essa granularidade alimenta o mapeamento direto entre comportamento observado e técnicas catalogadas pelo MITRE ATT\&CK, discutido na Seção~\ref{sec:lolbinattack}.

\subsection{A lacuna de telemetria do EID~1 e o \textit{health-gate}}
\label{subsec:eid1gap}

Um aspecto metodologicamente delicado, documentado honestamente nesta dissertação, é a \textbf{lacuna de telemetria do EID~1}. Na prática do laboratório, a esmagadora maioria dos documentos Sysmon corresponde a eventos de registro (EID~12/13), de modo que, em execuções de baixo sinal, as criações de processo dos LOLBin maliciosos podem ser esparsamente capturadas — ou ausentes — caso os agentes e o Sysmon não tenham concluído o \textit{check-in} junto ao Fleet no instante da execução do \textit{playbook}. Trata-se de uma limitação estrutural da instrumentação, não de uma falha de detecção: nenhum classificador pode reconhecer um processo cuja criação jamais foi registrada. A mitigação adotada é dupla. Primeiro, um \textbf{\textit{health-gate} de telemetria} substitui esperas fixas por um sinal de prontidão fim-a-fim: o ambiente só executa os \textit{playbooks} após confirmar, via Elasticsearch, que cada hospedeiro embarcou um número mínimo de eventos EID~1 frescos e que seu canal de execução remota está operante, sondando a cadeia completa \mbox{execução} $\rightarrow$ Sysmon $\rightarrow$ Fleet $\rightarrow$ Elasticsearch. Segundo, uma extração de janela priorizada por EID~1 garante a inclusão de todas as criações de processo disponíveis na janela de cada \textit{playbook}, antes de preencher o restante com eventos cronológicos, recuperando integralmente a cadeia de ataque em execuções de sinal rico. Essas salvaguardas são parte indissociável do \textit{pipeline} de medição e condicionam a validade interna das comparações entre as condições experimentais.

%======================================================================
\section{Técnicas LOLBin e o MITRE ATT\&CK}
\label{sec:lolbinattack}

\subsection{O framework MITRE ATT\&CK}
\label{subsec:attack}

O MITRE ATT\&CK é uma base de conhecimento, reconhecida internacionalmente, de táticas, técnicas e procedimentos (TTPs) adversariais observados no mundo real, organizada em táticas (o ``porquê'', como Persistência, Evasão de Defesa, Acesso a Credenciais, Coleta, Comando-e-Controle e Exfiltração) e técnicas (o ``como''). O framework oferece uma linguagem comum para descrever comportamento de atacante e mapear sistematicamente cada detecção a um identificador de técnica, o que confere rigor e interoperabilidade à caça. No MAS-Hunt, a lógica de detecção é explicitamente ancorada nas técnicas ATT\&CK relevantes — por exemplo, chaves de execução de registro (T1547.001), criação ou modificação de serviços do Windows (T1543.003), limpeza de logs de eventos (T1070.001), despejo de credenciais do sistema operacional (T1003), \textit{keylogging} (T1056.001), captura de tela (T1113) e exfiltração por canal de Comando-e-Controle (T1041).

\subsection{Binários do tipo Living-off-the-Land}
\label{subsec:lolbin}

As técnicas \textit{Living-off-the-Land} (LOLBin) consistem no abuso de binários, scripts e bibliotecas legítimos, já presentes no sistema operacional, para executar ações maliciosas. Por se valerem de ferramentas confiáveis e assinadas pelo próprio fornecedor — como \texttt{powershell.exe}, \texttt{certutil.exe}, \texttt{bitsadmin.exe}, \texttt{rundll32.exe} e \texttt{regsvr32.exe} —, essas técnicas evadem controles baseados em assinatura e em listas de permissão, dissolvendo a fronteira entre atividade administrativa benigna e atividade maliciosa. Justamente por isso, a discriminação entre uso legítimo e uso abusivo do mesmo binário exige análise comportamental contextual: o sinal não está na presença do executável, mas no encadeamento de seus argumentos, em sua filiação de processo e em sua correlação com conexões de rede e modificações de persistência. Essa dificuldade fundamenta o desenho do corpus experimental desta dissertação — \textbf{185 \textit{playbooks}} distribuídos por \textbf{nove famílias de LOLBin}, com instâncias maliciosas e benignas pareadas por família, de modo a forçar o sistema a decidir pelo comportamento e não pelo mero nome do binário. A predominância de exemplares benignos torna o cenário fiel à falácia da taxa-base (Seção~\ref{subsec:taxabase}) e justifica a adoção da métrica \textbf{F2}, que pondera a revocação (\textit{recall}) com peso quatro vezes superior ao da precisão ($\beta = 2$), refletindo o custo assimétrico de não detectar uma ameaça real frente ao de gerar um alarme adicional. À métrica F2 somam-se Precisão, Revocação e a Taxa de Falsos Positivos, todas reportadas com intervalos de confiança estimados por \textit{bootstrap}.

%======================================================================
\section{Injeção Adversarial de Logs}
\label{sec:injecaologs}

A injeção adversarial de logs é o vetor de ataque específico contra o qual a resiliência do MAS-Hunt é mensurada e que materializa, no domínio da caça a ameaças, a injeção de prompt indireta da Seção~\ref{subsec:injecao}. A premissa é que um adversário com a capacidade de escrever em logs — diretamente ou induzindo o sistema a registrar entradas controladas — pode incorporar à telemetria conteúdo textual cujo propósito não é descrever um evento, mas manipular o agente que virá a lê-lo. Um \textit{payload} de injeção pode, por exemplo, apresentar-se como uma diretiva de sistema instruindo o agente a ignorar a análise anterior e a classificar uma janela maliciosa como benigna, induzindo um falso negativo, ou o inverso, induzindo um falso positivo que sobrecarrega o analista. O perigo desse vetor é proporcional à autonomia do agente: quanto mais o agente confia na telemetria como fonte de verdade, mais exposto fica à telemetria adulterada.

O desenho experimental desta dissertação operacionaliza a ameaça em duas dimensões ortogonais. A primeira é o \textbf{modo de injeção}, que governa o momento e a relação temporal entre a injeção e a atividade-alvo — variando de rajadas prévias à execução (\textit{pre-burst}), passando por rajadas acopladas a cada ataque (\textit{per-attack-burst}), até a intercalação contínua (\textit{interleaved}) — além de um modo de controle sem injeção. A segunda é o \textbf{nível de furtividade} do \textit{payload}: \textit{overt} (instrução explícita em texto claro), \textit{encoded} (codificada, por exemplo em \textit{base64}) e \textit{steganographic} (oculta em campos aparentemente legítimos). A resiliência é então quantificada por meio do \textit{Attack Success Rate} (ASR) — a fração de injeções que efetivamente induziu a classificação incorreta — e do \textit{Non-Attributed Success Rate} (NASR), que restringe o sucesso aos casos em que o ataque não foi sequer detectado pelo agente, capturando, assim, não só se a defesa falhou, mas se falhou silenciosamente. A hipótese central da pesquisa é que a governança multiagente endurecida (C1-full) reduz, de forma estatisticamente significativa, o ASR e o NASR frente às linhas de base sem endurecimento (C1-nohardening e C2-naive), evidenciando que os mecanismos M1–M5 — em especial o M3 — conferem ao sistema resistência efetiva à injeção adversarial de logs.

%======================================================================
\section{Síntese do Capítulo}
\label{sec:sintese}

Este capítulo percorreu o caminho conceitual que vai da prática da caça a ameaças, motivada pela falência da detecção reativa e pela urgência de reduzir o tempo de permanência do atacante, até a ameaça específica da injeção adversarial de logs. Estabeleceu-se que os LLMs e os agentes de IA oferecem um mecanismo poderoso de automação do ciclo de caça dirigido por hipóteses, mas que sua autonomia e sua dependência do contexto textual os tornam vulneráveis a uma classe inédita de ataques — envenenamento de memória, injeção de prompt e amplificação de mau funcionamento. Argumentou-se que a orquestração multiagente hierárquica, fundamentada em evidência empírica de robustez sob adversidade, constitui resposta arquitetural a esses riscos, e que os cinco mecanismos de governança do MAS-Hunt (M1–M5) endereçam, cada um, uma face específica da ameaça. Por fim, fundamentaram-se a infraestrutura de telemetria Windows/Sysmon/Elastic sobre a qual o sistema opera, a lacuna de telemetria do EID~1 e sua mitigação por \textit{health-gate}, as técnicas LOLBin catalogadas pelo MITRE ATT\&CK que compõem o corpus experimental, e a racionalidade das métricas (F2, FPR, ASR, NASR) e das quatro condições (C1-full, C1-nohardening, C2-naive, C3-Elastic). Os capítulos seguintes detalham, respectivamente, a arquitetura do MAS-Hunt (Capítulo~\ref{cap:arquitetura}) e a metodologia experimental (Capítulo~\ref{cap:metodologia}) que põem à prova as proposições aqui fundamentadas.
